\documentclass[11pt]{article}

\usepackage[utf8]{inputenc}
\usepackage[T1]{fontenc}
\usepackage[english]{babel}
\usepackage{amsmath,amssymb,amsthm}
\usepackage{enumitem}
\usepackage{booktabs}
\usepackage{array}
\usepackage[hidelinks]{hyperref}
\usepackage{geometry}
\geometry{a4paper,margin=2.6cm}

\theoremstyle{plain}
\newtheorem{theorem}{Theorem}
\theoremstyle{remark}
\newtheorem*{remark}{Remark}

\newcommand{\Z}{\mathbb{Z}}
\newcommand{\N}{\mathbb{N}}
\newcommand{\R}{\mathbb{R}}
\newcommand{\Zi}{\mathbb{Z}[i]}
\newcommand{\pt}{\tilde{p}}
\newcommand{\lmult}{\lambda_{\mathrm{multiset}}}
\newcommand{\lset}{\lambda_{\mathrm{set}}}

\title{Prior Art for the Minimal-Period Formulas of\\
Rational Cut-and-Project Strip Projections:\\[2pt]
\large A Survey of the Search and the Remaining Verification Items}

\author{Dirk Kunert\thanks{The literature search summarised here was carried out
with the assistance of AI tools (ChatGPT, Gemini, and Claude). All bibliographic
data and exclusion judgements were checked against primary sources by the author;
remaining errors are the author's own.}}

\date{31 May 2026}

\begin{document}
\maketitle

\begin{abstract}
This note documents, for the benefit of a referee or arXiv endorser, the prior-art
search conducted in support of the manuscript on period lengths of one-dimensional
\emph{rational} cut-and-project sequences. The manuscript proves a closed-form
minimal period for the order-sorted gap sequence of a rational strip projection:
with coprime $\alpha,\beta\in\N$, window $\omega\ge 0$,
$N=\lfloor\omega\alpha\rfloor+\lfloor\omega\beta\rfloor+1$ and
$D=\alpha^2+\beta^2$, the multiset period equals $N$ when $D\nmid N$ and $N/D$
otherwise, while the set-valued period equals $N$ when $N<D$ and $1$ otherwise.
The concern motivating the search is that this result, or an equivalent lemma,
might already exist---possibly buried in work on an ostensibly different subject,
since the constant $D=\alpha^2+\beta^2$ is number-theoretically ubiquitous and the
ambient objects (strip projections, arithmetic discrete lines, gap sequences) are
classical. We argue that the \emph{conjunction} of features---a sorted projected
gap period, with $D=\alpha^2+\beta^2$ in the role of a period modulus, a
multiset-versus-set dichotomy, and a window dependence through $N$---forms a
distinctive fingerprint, and we report the outcome of four search rounds
(approximately $180$ logged queries across arXiv, HAL, dblp, Google Books, the
digital-geometry and aperiodic-order monograph literature, and the zbMATH Open
index), followed by individual treatment of six printed or poorly indexed sources.
No source examined states the result, as a theorem or as an auxiliary lemma. We
close by delineating the residual blind spots---principally the full-text bodies of
a few printed sources and the inherent limit of metadata- and review-level
indexing---and identify a single targeted expert enquiry as the most economical way
to discharge them.
\end{abstract}

\section{Purpose and scope}\label{sec:purpose}

The manuscript under consideration establishes minimal-period formulas for a
one-dimensional rational cut-and-project construction. Because its central constant
$D=\alpha^2+\beta^2$ occurs throughout number theory and discrete geometry, and
because the surrounding objects are classical, a natural concern---raised while
seeking an endorsement---is that the result may already be in the literature,
perhaps as an auxiliary lemma in a paper on a different topic. This note records the
resulting due-diligence search in a form an expert reader can audit: the axes
searched, the nearest catalogued neighbours and why each is structurally distinct,
and the residual uncertainty that no database search can, in principle, remove.

The document is deliberately self-contained and assumes familiarity with
cut-and-project (model) sets, the Steinhaus three-distance theorem, Sturmian and
Christoffel words, and the arithmetic-discrete-line tradition of digital geometry.
It is \emph{not} a contribution to the mathematics of the manuscript; it is a prior-art
statement. Section~\ref{sec:result} fixes the object and the fingerprint;
Section~\ref{sec:disguise} explains why the result is easy to overlook;
Section~\ref{sec:method} states the methodology;
Section~\ref{sec:findings} reports the findings axis by axis;
Section~\ref{sec:zbmath} reports the zbMATH Open round and how to read its empty
fingerprint queries; Section~\ref{sec:residual} lists what remains to be verified;
Section~\ref{sec:conclusion} states the balance. Appendix~\ref{app:queries}
reproduces the zbMATH query log.

\section{The result under assessment}\label{sec:result}

Fix coprime $\alpha,\beta\in\N$ and a window parameter $\omega\ge 0$. A lattice
point $(x,y)\in\Z^2$ is \emph{accepted} if it lies in the closed strip
\begin{equation}\label{eq:strip}
  -\beta\omega \;\le\; \alpha x-\beta y \;\le\; \alpha\omega ,
\end{equation}
and each accepted point is assigned the projected coordinate
$\pt=\beta x+\alpha y$, the image of $(x,y)$ under orthogonal projection onto the
line of slope $\alpha/\beta$ (up to the scaling that makes $\pt$ an integer). The
integer level $r:=\alpha x-\beta y$ ranges over the
\begin{equation}\label{eq:N}
  N=\lfloor\omega\alpha\rfloor+\lfloor\omega\beta\rfloor+1
\end{equation}
admissible values permitted by the closed strip~\eqref{eq:strip}; the constant term
$+1$ records the closed-boundary convention (lattice points on either bounding line
are accepted). Set
\begin{equation}\label{eq:D}
  D:=\alpha^2+\beta^2 .
\end{equation}
Because $\gcd(\alpha,\beta)=1$ we have $\gcd(\beta,D)=\gcd(\beta,\alpha^2)=1$ and
likewise $\gcd(\alpha,D)=1$, so $\beta$ is invertible modulo $D$ and the level and
projection coordinates of an accepted point satisfy
\begin{equation}\label{eq:residue}
  \pt \;\equiv\; -\alpha\,\beta^{-1}\,r \pmod{D}.
\end{equation}
Consequently the accepted values of $\pt$, taken in increasing order, form a
finite union of $N$ arithmetic progressions of common difference $D$, one per level
$r$; their consecutive differences (the \emph{gaps}) repeat periodically. Two
conventions are studied:
\begin{itemize}[leftmargin=1.8em,itemsep=2pt]
  \item the \emph{multiset} gap sequence, which retains the multiplicity of each gap
        value (distinct accepted points may share a projected coordinate);
  \item the \emph{set} gap sequence, which records only the distinct gap values in
        order.
\end{itemize}
The manuscript determines the minimal period of each.

\begin{theorem}[Multiset period]\label{thm:multiset}
\[
  \lmult \;=\;
  \begin{cases} N, & D\nmid N,\\[2pt] N/D, & D\mid N. \end{cases}
\]
\end{theorem}

\begin{theorem}[Set-valued period]\label{thm:set}
\[
  \lset \;=\;
  \begin{cases} N, & N<D,\\[2pt] 1, & N\ge D, \end{cases}
\]
and $\lset \mid \lmult$, with equality if and only if $N\le D$.
\end{theorem}

For an irrational slope the projected gap sequence admits no period at all; the
manuscript proves this aperiodicity separately and formalises it. Both period
theorems and the irrational case are machine-checked in Lean~4 against
\texttt{mathlib}.

\paragraph{The fingerprint.}
The object to be searched for is the \emph{conjunction} of four features:
\begin{enumerate}[label=(\roman*),leftmargin=2.2em,itemsep=2pt]
  \item the gap sequence of the orthogonally projected and \emph{order-sorted}
        accepted points---not the chain-code word read along the abscissa, and not a
        spectral or diffraction gap;
  \item the modulus $D=\alpha^2+\beta^2$ in the role of a \emph{period}
        (equivalently $N/D$ as a period in the degenerate case $D\mid N$);
  \item the \emph{multiset-versus-set} dichotomy, with the divisibility relation
        $\lset\mid\lmult$;
  \item the \emph{window dependence} through $N$ of~\eqref{eq:N}.
\end{enumerate}
Each feature in isolation is classical (see Section~\ref{sec:findings}); the claim
of novelty concerns their conjunction. The search below was designed to detect that
conjunction even when it is named differently or appears as a step in a larger
argument.

\section{Why the result is easy to overlook}\label{sec:disguise}

Two mathematical ``disguises'' make a direct keyword search inadequate and explain
where a hidden lemma would most plausibly live.

\paragraph{The strip is an arithmetic discrete line.}
Up to the choice of boundary convention, the strip~\eqref{eq:strip} coincides with
the arithmetic discrete line of Reveill\`es~\cite{Reveilles1991},
\[
  D(\alpha,\beta,\mu,W)=\{(x,y)\in\Z^2 : \mu \le \alpha x-\beta y < \mu+W\},
\]
with offset $\mu=-\beta\omega$ and arithmetic thickness $W=(\alpha+\beta)\omega$.
In digital geometry the same set carries an extensive and largely self-contained
vocabulary (naive, standard, and thick lines; the leaning/main direction vector
$(\beta,\alpha)$; remainder, period, and pattern), so a result about it could be
stated without ever mentioning ``cut and project''. The digital-geometry corpus is
therefore the single most likely hiding place, and it received the most attention
(Round~3, Section~\ref{ssec:dg}).

\paragraph{$\alpha^2+\beta^2$ is everywhere.}
The constant $D=\alpha^2+\beta^2$ is simultaneously the norm $N_{\Zi}(\beta+\alpha i)$
of a Gaussian integer, the squared length $\lVert(\beta,\alpha)\rVert^2$ of the
projection direction, and the determinant of the integer similarity aligning $\Z^2$
with the strip. It is thus number-theoretically natural and recurs across the
literature---in the two-squares theorems of Fermat and Euler, in Gaussian-integer
norms, in coincidence-site-lattice and similar-sublattice counting, and as a squared
radius in digital circles. The search must therefore distinguish the many
appearances of $\alpha^2+\beta^2$ \emph{as a counting or curvature quantity} from
its appearance \emph{as the period modulus} of the sorted projected gap sequence,
which is the feature claimed to be new. Every nearest neighbour identified below
turns out to use $\alpha^2+\beta^2$ in one of the former roles.

\section{Methodology and sources}\label{sec:method}

The search proceeded in four rounds and was logged for reproducibility: for each
query we recorded the date, the database, the exact query string, the hit count,
and the exclusion reason (most commonly: a different period definition; merely
qualitative periodicity; a symbolic word rather than the projected gap sequence;
no multiplicities; or a different notion of window).

\begin{itemize}[leftmargin=1.8em,itemsep=3pt]
  \item \textbf{Round~1} (notation- and topic-oriented; arXiv and web;
        $\approx 20$ queries) and \textbf{Round~2} (notation-independent, by
        mathematical substance: two coprime integers, the quantity $p^2+q^2$, a
        floor counting formula, multiset/set conventions; $\approx 28$ queries)
        swept the cut-and-project, three-gap, Sturmian, and quasicrystal literature.
  \item \textbf{Round~3} ($\approx 22$ queries; arXiv, HAL, dblp, Google Books, the
        AMS CRM series) targeted the arithmetic-discrete-line tradition identified in
        Round~2 as the most likely hiding place.
  \item Six printed or poorly indexed sources flagged after Round~3 were then worked
        individually (Section~\ref{ssec:books}).
  \item \textbf{Round~4} ($\approx 90$ field queries plus $24$ review-level
        fingerprint combinations) used the zbMATH Open REST API
        (\url{https://api.zbmath.org/v1/document/_search}) to make the search path
        exactly reproducible (Section~\ref{sec:zbmath}, Appendix~\ref{app:queries}).
\end{itemize}

Table~\ref{tab:sources} summarises the databases and their indexing depth. The
distinction between \emph{metadata/review-level} indexing (title, authors, MSC,
keywords, and---for zbMATH---the reviewer's summary text) and \emph{full-text}
indexing (the published PDF body) is essential to reading the result and is revisited
in Section~\ref{sec:residual}.

\begin{table}[h]
\centering
\small
\begin{tabular}{@{}lll@{}}
\toprule
Source & Indexing depth & Role in the search\\
\midrule
arXiv & title / abstract / full text & primary, post-2000\\
zbMATH Open & title / MSC / keywords / review text & primary, broad and historical\\
dblp & title / authors & enumeration of DGCI/DGMM proceedings\\
HAL & full text (theses) & Reveill\`es thesis; French DG school\\
Google Books & in-book search & monographs (Klette--Rosenfeld; Senechal)\\
NASA ADS / journal archives & abstract / full text & physics ``rational approximants''\\
\bottomrule
\end{tabular}
\caption{Databases consulted and the depth at which each is indexed.}
\label{tab:sources}
\end{table}

\section{Findings by literature axis}\label{sec:findings}

Across all axes the pattern is consistent: the constituent objects and constants are
classical and are cited in the manuscript as related work, but the conjunction of
Section~\ref{sec:result} was not located. We summarise each axis and then the
nearest individual records.

\subsection{Cut-and-project sets and model sets}

The model-set literature \cite{Meyer1972,BaakeGrimm2013,Senechal1995} is framed
throughout for \emph{totally irrational} windows, where the projected set is
aperiodic; the rational case is the trivial periodic crystal (a lattice Dirac comb)
and is generally not studied as a source of a nontrivial closed-form period. Gap and
patch statistics in this setting concern \emph{frequencies} and limiting
distributions for irrational slopes, not a finite minimal period for a rational one.
The nearest catalogued record on the cut-and-project axis,
G\"ahler--Hunton--Kellendonk~\cite{GaehlerHuntonKellendonk2013}, does treat
\emph{rational} projection-method patterns, but computes their integral \v{C}ech
cohomology---an algebraic-topological invariant of the associated tiling space, not a
gap period. Thus the phrase ``rational projection method'' is occupied in the
literature, but for cohomology.

\subsection{Arithmetic discrete lines and digital geometry}\label{ssec:dg}

This is the closest comparison and received the most scrutiny. The literature
contains exactly two genuine period statements about $D(\alpha,\beta,\mu,W)$, both
classical:
\begin{enumerate}[label=(\arabic*),leftmargin=2.2em,itemsep=2pt]
  \item \textbf{Translation periodicity} of the discrete line along its direction
        vector $(\beta,\alpha)$~\cite{Reveilles1991}, embedded in the modular
        $\mathrm{SL}(2,\Z)$ structure of such lines;
  \item the \textbf{chain-code / Christoffel-word period}---the code length $q$, or
        equivalently the word length $p+q$---of the naive line
        \cite{BerstelEtAl2008,LachaudSaid2013,KletteRosenfeld2004}.
\end{enumerate}
Three structural differences separate these from Theorems~\ref{thm:multiset}
and~\ref{thm:set}:
\begin{itemize}[leftmargin=1.8em,itemsep=2pt]
  \item \emph{Combinatorial object.} Digital geometry studies the chain-code word
        read in geometric order along $x$; the manuscript studies the
        \emph{order-sorted} sequence of projected distances. These are different
        sequences with different periods.
  \item \emph{Window dependence.} The digital-geometry period depends only on the
        slope; the manuscript's $N$ depends centrally on the window $\omega$.
  \item \emph{Multiset convention.} Discrete lines are point sets; no multiset
        (multiplicity-bearing) convention, and hence no analogue of
        $\lset\mid\lmult$, appears in that literature.
\end{itemize}
Where $\alpha^2+\beta^2$ does occur in digital geometry it is a squared radius or
curvature for discrete circles
\cite{FiorioJametToutant2006}---not a period modulus. Three specific verifications
were carried out:
\begin{itemize}[leftmargin=1.8em,itemsep=2pt]
  \item The Reveill\`es thesis~\cite{Reveilles1991} was read in full from the HAL
        scan ($245$\,pp.). Its periodicity results are exactly (1) and (2) above; the
        thematically closest section, \S VIII (``Lattice problem''; \S VIII.3
        ``Droites discr\`etes et r\'esidus quadratiques''), treats Gauss-sum
        lattice-point counting, with \emph{no} sorted gap sequence, \emph{no}
        multiset/set distinction, and $\alpha^2+\beta^2$ \emph{not} as a period
        modulus.
  \item The DGCI/DGMM proceedings (all $20$ dblp-indexed volumes, 1996--2025) were
        title-scanned for \texttt{period}, \texttt{gap}, \texttt{distance},
        \texttt{multiset}, \texttt{Christoffel}, \texttt{Sturm}, \texttt{pattern},
        \texttt{chain code}, \texttt{thickness}, \texttt{residue},
        \texttt{projection}, \texttt{lattice}; the four nearest candidates were then
        checked at full text. All are recognition algorithms or different objects
        (Sivignon, Farey-fan subsegment recognition; Laboureix--Debled-Rennesson,
        Stern--Brocot recognition; Roussillon--Labb\'e, rational discrete planes;
        Hoarau--Monteil, patterns in discrete circles via Pell--Fermat); none
        concerns a sorted projected gap period.
  \item Klette--Rosenfeld, \emph{Digital Geometry}~\cite{KletteRosenfeld2004},
        Ch.~9 (``2D Straightness'') was verified by in-book search: the word
        ``multiset'' occurs nowhere in the book, and ``orthogonal projection''
        appears only in Ch.~8 (arc length and surface area). The companion review
        article characterises rational-DSS periodicity as ``the points of any
        rational DSS are defined by a basic set of points''---the slope-bound
        chain-code/pattern period, exactly line~(2) above.
\end{itemize}

\subsection{The three-distance theorem; Sturmian and Christoffel words}

The Steinhaus three-distance theorem yields the qualitative statement that at most
three distinct gaps occur~\cite{ThreeGap}; it gives no period \emph{length}.
Sturmian and Christoffel theory consistently yields the period $p+q$, the word
length of the characteristic sequence~\cite{BerstelEtAl2008}---linear in $p,q$,
never $p^2+q^2$. These are the canonical $p+q$ phenomena; the manuscript cites them
as related, not competing, results.

\subsection{Monographs and festschriften on aperiodic order}\label{ssec:books}

Baake--Grimm, \emph{Aperiodic Order}, Vols.~1 and~2
\cite{BaakeGrimm2013,BaakeGrimm2017}, and Senechal, \emph{Quasicrystals and
Geometry}~\cite{Senechal1995}, are framed uniformly for irrational/aperiodic
windows; rational slopes appear only as the trivial ideal crystal. Senechal was
verified at searchable full text ($306$\,pp.): ``multiset'' does not occur, ``gap''
appears only in the sense of gapless tilings, and the only sum-of-squares norm is
the four-dimensional quaternionic norm of the $E_8$/icosahedral construction. The
sole $\alpha^2+\beta^2$-adjacent appearance in either monograph is in
coincidence-site-lattice / similar-sublattice counting (Baake--Zeiner, Vol.~2,
Ch.~3), where $a^2+b^2$ is a Gaussian-integer norm in a Dirichlet series---again a
counting constant, not a period. In the festschrift \emph{Directions in Mathematical
Quasicrystals}~\cite{BaakeMoody2000}, the contribution flagged as most suspect,
Pleasants' ``Designer quasicrystals''~\cite{Pleasants2000}, constructs
\emph{aperiodic} cut-and-project sets with pre-assigned patch and diffraction
properties; ``gap'' in that volume means a spectral gap (gap labelling), not a gap
sequence.

\subsection{Rational approximants in the physics literature}

In physics, a ``rational approximant'' replaces the golden-mean slope $\tau$ by its
convergents $\tfrac11,\tfrac21,\tfrac32,\tfrac53,\dots$; the $k$-th approximant's
unit-cell period is a Fibonacci number, linear in the order---again the $p+q$
phenomenon, not $p^2+q^2$. Korepin's integrable quasicrystal models live on Penrose
tilings and are unrelated to a one-dimensional gap period. This axis is therefore
\emph{structurally} excluded rather than exhaustively searched; see
Section~\ref{sec:residual}.

\subsection{Triage of the nearest catalogued records}

Table~\ref{tab:triage} lists the four catalogued records that came closest on
vocabulary, each with the reason it is not the result.

\begin{table}[h]
\centering
\small
\renewcommand{\arraystretch}{1.25}
\begin{tabular}{@{}>{\raggedright}p{4.6cm}>{\raggedright}p{2.4cm}p{7.3cm}@{}}
\toprule
Record & Identifier & Why it is not the result\\
\midrule
Haynes--Koivusalo--Walton--Sadun, ``Gaps problems and frequencies of patches in cut and project sets'' \cite{HaynesEtAl2016}
& Zbl 1371.11110 (2016)
& Right authors and vocabulary (``gaps'' $+$ ``cut and project''), but patch-frequency statistics of \emph{irrational} model sets; no closed minimal period, no $\alpha^2+\beta^2$ modulus.\\
G\"ahler--Hunton--Kellendonk, ``Integral cohomology of rational projection method patterns'' \cite{GaehlerHuntonKellendonk2013}
& Zbl 1270.52025 (2013)
& ``Rational'' $+$ ``projection,'' but a topological invariant (\v{C}ech cohomology of tiling spaces); not a gap period.\\
Cotfas, ``Modified strip-projection method'' \cite{Cotfas2006}
& arXiv:math-ph/0605046
& Nearest vocabulary match (``strip projection''); occupation density of quasiperiodic packings, irrational slope, no period.\\
M\"uller--Dressler, ``\"Uber eine gewichtete Mittelung der Gitterpunkte'' \cite{MuellerDressler1972}
& Zbl 0227.10040 (1972)
& Only hit of the ``sum of two squares'' $+$ ``lattice'' $+$ ``period'' query; weighted lattice-point averaging, not a projected gap period.\\
\bottomrule
\end{tabular}
\caption{The nearest catalogued records and the reason each is excluded.}
\label{tab:triage}
\end{table}

\section{The zbMATH Open round}\label{sec:zbmath}

Round~4 queried the zbMATH Open index programmatically through its REST API, so that
the search path is exactly reproducible. Three properties of the API matter for
interpreting the result.

\begin{enumerate}[leftmargin=2.0em,itemsep=3pt]
  \item The field \texttt{ab:} does not search a separate abstract but the
        \emph{reviewer's summary text} (the \texttt{editorial\_contributions}
        field); \texttt{any:} includes it. This text is indexed and searchable even
        when its \emph{display} is licence-gated (some records return the placeholder
        ``zbMATH Open Web Interface contents unavailable due to conflicting
        licenses,'' while others, e.g.\ the review of
        \emph{Aperiodic Order} Vol.~1, return the full text). Searchability is
        unaffected by display gating: one searches the index, not the rendered page.
  \item A review-/summary-level search is therefore possible, not merely a
        title/MSC search (calibration: \texttt{ab:period} returns $52\,753$ hits,
        \texttt{any:period} returns $60\,534$). What is \emph{not} searched is the
        published PDF body.
  \item The hit count is returned in \texttt{status.nr\_total\_results}; an
        HTTP~404 response denotes \emph{zero} hits (body: ``successful access, but no
        result''), not a syntax error.
\end{enumerate}

Calibration queries confirm the runs are intact (\texttt{ti:group} $=60\,517$,
\texttt{au:Serre} $=657$). The decisive observations are the \emph{empty}
fingerprint queries at the review/all-fields level---for example
\texttt{any:"period length" \& any:"cut and project"} $=0$,
\texttt{any:"minimal period" \& any:projection \& any:gap} $=0$, and
\texttt{ab:"gap sequence" \& ab:period} $=0$. The conjoined invariant surfaces at
no title-, review-, or keyword-level of the zbMATH index. The complete diagnostic
log is reproduced in Appendix~\ref{app:queries}.

\section{Residual verification items}\label{sec:residual}

The search lowers the residual risk to \emph{low}, but cannot reduce it to zero. We
state the remaining items honestly, ordered by importance.

\begin{enumerate}[leftmargin=2.0em,itemsep=4pt]
  \item \textbf{Full-text (PDF-body) blind spot --- the principal, irreducible
        limit.} Metadata- and review-level indexing cannot detect a result used as
        an auxiliary lemma in a paper whose \emph{title and reviewer summary both}
        avoid the words period, gap, and projection. This is a structural limit of
        every database consulted. Mitigant: the fingerprint---a sum of two squares
        as the period modulus of a \emph{sorted} gap sequence, with a multiset/set
        dichotomy---is linguistically distinctive, and would most likely surface as
        a keyword or summary term if genuinely treated; this is probable, not
        guaranteed.
  \item \textbf{Baake--Grimm Vol.~1, verbatim index check.} Only the frontmatter and
        corrigenda were available; the verbatim index entries ``rational'' and
        ``periodic case'' were not read. Residual risk is low, given the uniformly
        aperiodic framing of both volumes.
  \item \textbf{Pleasants' contribution, full text.} Only the volume's table of
        contents was inspected. The contribution is explicitly about aperiodic
        designer model sets; a verbatim reading of the printed article would close
        the item.
  \item \textbf{Pre-2000 physics.} Excluded structurally (approximant period
        $\sim p+q$, a Fibonacci number), but this literature is large and poorly
        web-indexed, so the exclusion is by argument rather than by exhaustive
        search.
  \item \textbf{MathSciNet and NASA ADS not queried programmatically.} Their
        coverage overlaps zbMATH and arXiv; the marginal value lies only in
        item~4 above.
\end{enumerate}

\paragraph{Recommended next step.}
The most economical way to discharge items~1--4 is \emph{not} a further database run
but a single targeted enquiry to an expert in arithmetic discrete lines or aperiodic
order, who can confirm from working knowledge of the full-text corpus whether the
conjoined invariant of Section~\ref{sec:result} is novel. A database search has been
pushed to its natural limit; the remaining uncertainty is exactly the kind a
specialist can resolve from memory.

\section{Balance}\label{sec:conclusion}

After approximately $180$ logged queries across arXiv, the web, HAL, dblp, Google
Books, and the zbMATH Open index, together with individual treatment of six printed
or poorly indexed sources, no source examined states the conjoined result of
Theorems~\ref{thm:multiset} and~\ref{thm:set}, as a theorem or as an auxiliary
lemma. The constituent constants and objects---$\alpha^2+\beta^2$ as a Gaussian norm
or curvature, the translation and chain-code periodicities of arithmetic discrete
lines, the three-distance bound, the Sturmian/Christoffel period $p+q$, the rational
crystal as a degenerate model set---are classical and are cited as such. What was
not located is their conjunction: a closed-form minimal period of the order-sorted
projected gap sequence of a rational strip projection, with modulus
$D=\alpha^2+\beta^2$, a multiset/set dichotomy, and a window dependence through $N$.

The residual risk is low and is bounded explicitly by the full-text blind spot of
Section~\ref{sec:residual}. The result is therefore presented as, to the best of a
thorough but necessarily incomplete search, new; the author welcomes correction from
any reader who recognises it.

\appendix
\section{Reproducible zbMATH query log}\label{app:queries}

Queries use field prefixes (\texttt{ti:}, \texttt{au:}, \texttt{cc:}, \texttt{ab:},
\texttt{any:}), quoted phrases, and the conjunction operator `` \& ''. Counts are
\texttt{status.nr\_total\_results} as of 30~May 2026; ``$0$'' denotes an HTTP~404
(no result). Parenthetical notes give the dominant or only record behind a small
nonzero count.

\begin{small}
\begin{verbatim}
-- Calibration (runs intact) --
ti:group                                        60517
au:Serre                                          657
ab:period                                       52753
any:period                                      60534

-- Direct axis (title level) --
ti:"cut-and-project"                               48
ti:"cut-and-project" & ti:period                    0
ti:"cut-and-project" & ti:gap                       0
ti:"model set"       & ti:period                    0
cc:52C23 & ti:period                                0
cc:52C23 & ti:gap                                  10   (gap labelling / spectral gaps)
cc:52C23 & ti:rational                              3   (cohomology / tilings)
ti:"strip projection"                               1   (Cotfas 2006, irrational)

-- Digital geometry (title level) --
ti:"digital straightness"                           5
ti:"chain code"                                    38
au:Reveilles                                       16
ti:"discrete line" & ti:period                      0
ti:"arithmetic discrete line"                       0

-- Related concepts (title level) --
ti:"three distance" / ti:"three-gap"           22 / 18
ti:Sturmian & ti:period / & ti:gap               2 / 2
ti:Christoffel & ti:period                          1   (Schwarz-Christoffel)
ti:Beatty & ti:gap / & ti:period                 0 / 0

-- Fingerprint (review / all-fields level) --
any:"cut and project" & any:period                  2   (Baake-Grimm; off-topic)
ab:"cut and project" & ab:period                    1   (Baake-Grimm; off-topic)
any:"period length" & any:"cut and project"         0
any:"minimal period" & any:projection & any:gap     0
any:residue & any:projection & any:gap & any:period 0
ab:"gap sequence" & ab:period                       0
any:"sum of two squares" & any:period & any:lattice 1   (Mueller-Dressler 1972)
ab:"arithmetic discrete line" & ab:period           0
ab:"model set" & ab:period & ab:rational            0
\end{verbatim}
\end{small}

\begin{thebibliography}{99}

\bibitem{BaakeGrimm2013}
M.~Baake and U.~Grimm,
\emph{Aperiodic Order. Vol.~1: A Mathematical Invitation},
Encyclopedia of Mathematics and its Applications, vol.~149,
Cambridge University Press, 2013.

\bibitem{BaakeGrimm2017}
M.~Baake and U.~Grimm,
\emph{Aperiodic Order. Vol.~2: Crystallography and Almost Periodicity},
Encyclopedia of Mathematics and its Applications, vol.~166,
Cambridge University Press, 2017.

\bibitem{BaakeMoody2000}
M.~Baake and R.~V.~Moody (eds.),
\emph{Directions in Mathematical Quasicrystals},
CRM Monograph Series, vol.~13, American Mathematical Society, 2000.

\bibitem{BerstelEtAl2008}
J.~Berstel, A.~Lauve, C.~Reutenauer, and F.~V.~Saliola,
\emph{Combinatorics on Words: Christoffel Words and Repetitions in Words},
CRM Monograph Series, vol.~27, American Mathematical Society, 2008.

\bibitem{Cotfas2006}
N.~Cotfas,
``A modified strip-projection method,''
arXiv:math-ph/0605046, 2006.

\bibitem{FiorioJametToutant2006}
C.~Fiorio, D.~Jamet, and J.-L.~Toutant,
``Discrete circles: an arithmetical approach with non-constant thickness,''
in \emph{Vision Geometry XIV}, Proc.\ SPIE, vol.~6066, 2006.
HAL:~hal-00580556.

\bibitem{GaehlerHuntonKellendonk2013}
F.~G\"ahler, J.~Hunton, and J.~Kellendonk,
``Integral cohomology of rational projection method patterns,''
\emph{Algebraic \& Geometric Topology} \textbf{13}(3) (2013) 1661--1708.
\textsc{doi}:\,\href{https://doi.org/10.2140/agt.2013.13.1661}{10.2140/agt.2013.13.1661}.
Zbl~1270.52025.

\bibitem{HaynesEtAl2016}
A.~Haynes, H.~Koivusalo, J.~Walton, and L.~Sadun,
``Gaps problems and frequencies of patches in cut and project sets,''
\emph{Mathematical Proceedings of the Cambridge Philosophical Society}
\textbf{161}(1) (2016) 65--85.
Zbl~1371.11110.

\bibitem{KletteRosenfeld2004}
R.~Klette and A.~Rosenfeld,
\emph{Digital Geometry: Geometric Methods for Digital Picture Analysis},
Morgan Kaufmann, 2004;
companion review: ``Digital straightness---a review,''
\emph{Discrete Applied Mathematics} \textbf{139} (2004) 197--230.

\bibitem{LachaudSaid2013}
J.-O.~Lachaud and M.~Sa\"{\i}d,
``Two efficient algorithms for computing the characteristics of a subsegment
of a digital straight line,''
\emph{Discrete Applied Mathematics} \textbf{161}(15) (2013) 2293--2315.
\textsc{doi}:\,\href{https://doi.org/10.1016/j.dam.2012.08.038}{10.1016/j.dam.2012.08.038}.

\bibitem{Meyer1972}
Y.~Meyer,
\emph{Algebraic Numbers and Harmonic Analysis},
North-Holland, Amsterdam, 1972.

\bibitem{MuellerDressler1972}
W.~M\"uller and A.~Dressler,
``\"Uber eine gewichtete Mittelung der Gitterpunkte in der Ebene,''
1972.
Zbl~0227.10040.

\bibitem{Pleasants2000}
P.~A.~B.~Pleasants,
``Designer quasicrystals: cut-and-project sets with pre-assigned properties,''
in \cite{BaakeMoody2000}.

\bibitem{Reveilles1991}
J.-P.~Reveill\`es,
\emph{G\'eom\'etrie discr\`ete, calcul en nombres entiers et algorithmique},
Th\`ese d'\'Etat, Universit\'e Louis Pasteur, Strasbourg, 1991.
HAL:~tel-01279525.

\bibitem{Senechal1995}
M.~Senechal,
\emph{Quasicrystals and Geometry},
Cambridge University Press, 1995.

\bibitem{ThreeGap}
V.~T.~S\'os,
``On the distribution mod~$1$ of the sequence $n\alpha$,''
\emph{Ann.\ Univ.\ Sci.\ Budapest.\ E\"otv\"os Sect.\ Math.} \textbf{1} (1958) 127--134;
see also the three-distance (three-gap) theorem and its standard references.

\end{thebibliography}

\end{document}
